\documentclass[10pt]{article}
\usepackage{icmlko}

\icmlmeta{IP 파운데이션 월드모델}{138}{플래그 하나가 축 다섯을 이긴다}{바닥선을 무작위로 잡고 있었다 --- 제대로 잡으니 웹툰에서 ``완결인가'' 한 칸이 노트 5부터 쌓은 것을 이긴다}

\begin{document}

\twocolumn[
\icmlheader

\begin{abstract}
노트 137이 판 $\rho$의 사분의 일이 집단 간 순서라는 것을 보였다. 그러면
물어야 할 것이 하나 남는다 --- \textbf{집단만 아는 예측기는 얼마나 잘하나?}
포트폴리오의 바닥선이 \texttt{F0} 무작위였는데, 그건 바닥선이 아니다.
누구나 공짜로 아는 것이 바닥선이어야 한다. 재 봤다. \textbf{그리고 한 번
틀렸다} --- 처음엔 검정 집합 자신의 집단 평균으로 예측기를 만들어
세 도메인 판 $+$0.3589를 얻고 ``손 축 다섯이 플래그 셋보다 못하다''고
읽을 뻔했다. \textbf{그건 누수다.} 학습 집합의 집단 평균을 배워 검정에
적용하는 정직한 판은 $+$0.1255다. 그 차이가 어디서 왔는지도 재밌다 ---
웹툰과 모바일은 누수판과 정직판이 \emph{완전히 같고}(둘 다 $+$0.3195 ·
$+$0.4499), \textbf{애니만 $+$0.339에서 $-$0.338로 뒤집힌다.} 매체
(TVA · 극장판 · OVA)와 라벨의 관계가 2025년을 넘으며 반대가 된다 ---
노트 126의 법칙이 메타데이터 플래그에서 나타난 사례다. \textbf{정직한
비교의 결과는 이렇다.} 웹툰에서 ``완결인가'' 한 칸이 $+$0.319로 손 축
다섯 모형의 $+$0.202를 \textbf{이긴다.} 모바일에서 ``무료인가'' 한 칸이
$+$0.450으로 모형의 $+$0.448과 \emph{같다}. 애니에서만 모형이 크게 앞선다
($+$0.368 대 $-$0.338). 판으로 합치면 모형이 이기지만($+$0.321 대
$+$0.126) 그 우세는 애니 하나에서 온다. \textbf{그리고 외부 검색 축이
그 그림을 바꾼다} --- 세 도메인 모두에서 플래그를 넘어선다(웹툰 $+$0.354 ·
모바일 $+$0.551 · 애니 $+$0.487).
\end{abstract}
]

\begin{figure*}[t]
\centering
\includegraphics[width=\textwidth]{figs/floor.pdf}
\caption{왼쪽 --- 웹툰에서 플래그 하나가 축 다섯을 이긴다. 오른쪽 ---
바닥선도 누수될 수 있다. 애니만 부호가 뒤집힌다.}
\end{figure*}

\section{바닥선을 잘못 잡고 있었다}

포트폴리오의 바닥선이 \texttt{F0\_chance}(무작위 점수)다. 노트 125에서
넣을 때는 ``우연보다 나은가''를 재려던 것이었는데, 노트 137이 판 $\rho$의
사분의 일이 집단 간 순서라는 것을 보이고 나니 질문이 달라진다.
\textbf{누구나 공짜로 아는 것보다 나은가?}

세 도메인에 표지가 있다. 웹툰은 완결 여부, 모바일은 무료 여부, 애니는
매체다. 셋 다 기획 시점에 아는 것이고 축을 하나도 안 매겨도 알 수 있다.

\section{그리고 한 번 틀렸다}

처음 만든 예측기는 각 레코드를 \emph{그 집단의 라벨 평균 순위}로 예측했다.
세 도메인 판이 $+$0.3589가 나왔고, 손 축 다섯으로 만든 모형이 $+$0.3208
이니 ``노트 5부터 쌓은 것이 플래그 셋보다 못하다''로 읽힌다.

\claimx{\textbf{그 집단 평균을 검정 집합에서 뽑았다 --- 누수다.} 학습
집합에서 배워 검정에 적용해야 맞다. 정직한 판은 $+$0.1255다.}{바닥선을
만들 때도 시간 분할을 지켜야 한다는 것이 새삼스럽지 않은데, ``집단 평균''
이 모수처럼 안 보여서 놓쳤다. \textbf{외우는 것은 전부 모수다.}}

\begin{table}[t]
\centering\tsize
\caption{집단만 아는 예측기. 같은 표지, 같은 검정 집합.}
\begin{tabular}{lrrr}
\toprule
& $n$ & 누수 & 정직 \\
\midrule
웹툰 (완결 여부) & 711 & $+$.3195 & $+$.3195 \\
모바일 (무료 여부) & 441 & $+$.4499 & $+$.4499 \\
\textbf{애니 (매체)} & 606 & $+$.3390 & \textbf{$-$.3382} \\
\midrule
세 도메인 판 & 1{,}758 & $+$.3589 & $+$.1255 \\
\bottomrule
\end{tabular}
\end{table}

\claimx{\textbf{웹툰과 모바일은 누수판과 정직판이 완전히 같다.} 두 표지는
2025년을 넘어도 방향이 그대로다. \textbf{애니만 $+$0.339에서 $-$0.338로
뒤집힌다} --- 매체와 라벨의 관계가 반대가 됐다.}{애니 라벨은 라프텔 리뷰
수다. 2025년 이후 극장판과 TVA의 리뷰 수 관계가 왜 뒤집혔는지는 이 자료로
모른다. \textbf{노트 126의 법칙이 학습된 모수가 아니라 메타데이터
플래그에서 나타난 사례}라는 점이 새롭다.}

\section{정직하게 견주면}

\begin{table}[t]
\centering\tsize
\caption{같은 검정 집합. 판은 세 도메인 표본 가중.}
\begin{tabular}{lrrrr}
\toprule
& 웹툰 & 모바일 & 애니 & 판 \\
\midrule
집단만 (정직) & \textbf{$+$.319} & $+$.450 & $-$.338 & $+$.126 \\
F8 손 축 다섯 & $+$.202 & $+$.448 & $+$.368 & $+$.321 \\
F8 $+$검색 & $+$.347 & $+$.522 & $+$.430 & $+$.420 \\
F8 $+$검색$+$달력 & $+$.354 & \textbf{$+$.551} & \textbf{$+$.487} & \textbf{$+$.449} \\
\bottomrule
\end{tabular}
\end{table}

\claimx{\textbf{웹툰에서 플래그 하나가 축 다섯을 이긴다} --- $+$0.319 대
$+$0.202. 모바일에서는 비긴다 --- $+$0.450 대 $+$0.448. 노트 5부터 124까지
쌓은 축 다섯이, 두 큰 도메인에서 메타데이터 한 칸에 밀리거나 비긴다.}{판으로
합치면 모형이 이긴다($+$0.321 대 $+$0.126). 다만 그 우세는 \emph{애니
하나}에서 온다 --- 거기서만 플래그가 뒤집히기 때문이다. 플래그가 안정적인
두 도메인에서는 우세가 없다.}

\claimx{\textbf{외부 검색 축이 그 그림을 바꾼다.} 검색을 붙이면 세 도메인
모두에서 플래그를 넘어선다 --- 웹툰 $+$0.347, 모바일 $+$0.522, 애니
$+$0.430. 달력까지 붙이면 $+$0.354 · $+$0.551 · $+$0.487이다.}{노트 128이
검색 축을 ``작품 바깥에서 온 첫 축''이라 불렀는데, 이 표를 보면 그것이
\textbf{모형이 공짜 정보를 처음으로 확실히 넘어선 지점}이기도 하다.}

\section{바닥선을 갈아 끼웠다}

\texttt{F17\_grouponly}를 포트폴리오에 바닥선으로 올렸다. 집단 표지를
축으로 만들어(\texttt{lab/grpaxes.py}) 그 축 하나만 쓰는 정식화다. 여덟
도메인에 표지를 물렸다 --- 완결 여부 · 무료 여부 · 매체 · 포맷 · 국가 ·
카테고리 · 출판사.

\claimx{\textbf{바닥선은 무작위가 아니라 공짜로 아는 것이어야 한다.}
\texttt{F0\_chance}는 판 $+$0.009로 아무것도 안 걸러 낸다. 어떤 모형이든
그것보다는 낫다. 걸러 내야 할 것은 ``메타데이터 한 칸으로 되는 일을 축
스물세 개로 하고 있지 않은가''다.}{\texttt{F0} 는 남겨 둔다 --- 치환 귀무와
함께 ``우연보다 나은가''를 재는 데는 여전히 쓸모가 있다.}

\begin{table}[t]
\centering\tsize
\caption{남은 것.}
\begin{tabular}{ll}
\toprule
애니 & 매체 플래그가 왜 뒤집혔나 --- 라벨 쪽 변화 의심 \\
표지 & 나머지 도메인 표지의 정직판을 다 잴 것 \\
축 다섯 & 웹툰에서 플래그에 지는 이유를 갈라야 한다 \\
\bottomrule
\end{tabular}
\end{table}

마지막 줄이 다음 자리다. 웹툰에서 손 축 다섯이 $+$0.202이고 완결 플래그가
$+$0.319다. 축이 완결 여부를 못 담고 있다는 뜻인데, \textbf{그 정보는
축을 매길 때 눈앞에 있었다.} 축 정의가 무엇을 빠뜨렸는지 보는 것이
스물세 개를 더 붙이는 것보다 값질 수 있다.

\end{document}
